\chapter{Real Numbers}\label{ch:b1:reals}

All of analysis rests on one property that distinguishes $\R$ from
$\Q$: every nonempty set bounded from above has a \emph{least} upper
bound. This chapter states it precisely, derives its first
consequences --- the Archimedean property, the floor function, the
density of the rationals and of the irrationals --- and sets up the
vocabulary (sup, inf, max, min) used constantly from
\cref{ch:b1:seq} onward.

\section{The upper bound property}

\begin{definition}[Bounds, sup and inf]\label{def:b1:reals:bounds}
Let $A \subseteq \R$ be nonempty. A real $M$ is an \emph{upper
bound}\index{upper bound} of $A$ when $a \leq M$ for all $a \in A$;
$A$ is \emph{bounded above} when it has an upper bound (similarly
below, with lower bounds; \emph{bounded} means both). A
\emph{maximum} of $A$ is an upper bound belonging to $A$.

The \emph{supremum} $\sup A$\index{supremum} is the least upper bound
of $A$, when it exists; the \emph{infimum} $\inf A$\index{infimum} is
the greatest lower bound.
\end{definition}

\begin{theorem}[Completeness axiom of $\R$]\label{thm:b1:reals:sup}
$\R$ is an ordered field containing $\Q$ in which \emph{every
nonempty subset bounded above has a supremum}.\index{completeness of
R@completeness of $\R$}
\end{theorem}
\admitted

\begin{remark}
We take this as the defining axiom of $\R$; constructing a model (by
Dedekind cuts or by Cauchy sequences of rationals) and proving its
uniqueness is honest but long, and is left for further study. Note
that $\Q$ fails the property: $\{x \in \Q : x^2 < 2\}$ is bounded
above but has no least upper bound \emph{in $\Q$} --- its candidate,
$\sqrt 2$, is missing (\cref{ex:b1:logic:sqrt2}). By passing to
opposites ($\sup(-A) = -\inf A$), every nonempty set bounded below
has an infimum.
\end{remark}

\begin{proposition}[The $\varepsilon$-characterization]\label{prop:b1:reals:epsilon}
Let $A \neq \emptyset$ be bounded above and $s \in \R$. Then $s =
\sup A$ if and only if
\begin{enumerate}
  \item $s$ is an upper bound: $\forall a \in A$, $a \leq s$; and
  \item nothing smaller is: $\forall \varepsilon > 0$, $\exists a \in
        A$, $a > s - \varepsilon$.
\end{enumerate}
\end{proposition}

\begin{proof}
If $s = \sup A$: (1) holds by definition, and for (2), $s -
\varepsilon < s$ is not an upper bound, which is exactly the
existence of $a > s - \varepsilon$. Conversely, (1) says $s$ is an
upper bound; (2) says no $t < s$ is an upper bound (take $\varepsilon
= s - t$): $s$ is the least one.
\end{proof}

\begin{example}\label{ex:b1:reals:supexamples}
$\sup \intoo{0}{1} = 1$, not attained (no maximum);
$\sup \intcc{0}{1} = 1 = \max$. For $A = \{1 - \frac 1n : n \in
\N^*\}$: $\sup A = 1$, not attained; $\inf A = \min A = 0$. A maximum,
when it exists, is the supremum; the whole point of $\sup$ is to have
a substitute when the maximum does not exist.
\end{example}

\section{Consequences of completeness}

\begin{theorem}[Archimedean property]\label{thm:b1:reals:archimedes}
For every $x \in \R$ there is $n \in \N$ with $n > x$. Equivalently:
for all $\varepsilon > 0$ and $y > 0$, some multiple $n\varepsilon$
exceeds $y$.\index{Archimedean property}
\end{theorem}

\begin{proof}
Suppose not: some $x$ is an upper bound of $\N$. Then $s = \sup \N$
exists (\cref{thm:b1:reals:sup}). By
\cref{prop:b1:reals:epsilon}~(2) with $\varepsilon = 1$, there is $n
\in \N$ with $n > s - 1$; but then $n + 1 \in \N$ and $n + 1 > s$,
contradicting that $s$ is an upper bound. The second form follows
with $x = y/\varepsilon$.
\end{proof}

\begin{theorem}[Floor function]\label{thm:b1:reals:floor}
For every $x \in \R$ there is exactly one integer, the
\emph{floor}\index{floor function} $\lfloor x \rfloor$, with
\[
\lfloor x \rfloor \leq x < \lfloor x \rfloor + 1 .
\]
\end{theorem}

\begin{proof}
\emph{Existence.} The set $E = \{k \in \Z : k \leq x\}$ is nonempty:
by \cref{thm:b1:reals:archimedes} there is $m \in \N$ with $m > -x$,
and then $-m < x$, so $-m \in E$. It is bounded above (by any integer
$n > x$, which exists for the same reason), so, being a set of
integers trapped in the finite range $\intint{-m}{n}$, it has a
greatest element $k = \max E$. Then $k \leq x$, and $k + 1 \notin E$
means $x < k + 1$.

\emph{Uniqueness.} If $k$ and $k'$ both satisfy the inequalities,
then $k \leq x < k' + 1$ gives $k \leq k'$, and symmetrically $k'
\leq k$.
\end{proof}

\begin{theorem}[Density of $\Q$ and of $\R \setminus \Q$]\label{thm:b1:reals:density}
Between any two reals $x < y$ there lie a rational and an
irrational.\index{density}
\end{theorem}

\begin{proof}
\emph{A rational.} By \cref{thm:b1:reals:archimedes}, pick $n \in
\N^*$ with $n > \frac{1}{y - x}$, so $ny - nx > 1$. Let $m = \lfloor
nx \rfloor + 1$. On one side, $nx < \lfloor nx \rfloor + 1 = m$
(\cref{thm:b1:reals:floor}); on the other, $m = \lfloor nx \rfloor +
1 \leq nx + 1 < ny$. Dividing by $n$: $x < \frac mn < y$.

\emph{An irrational.} Apply the previous point to the pair $x -
\sqrt 2 < y - \sqrt 2$: some rational $q$ lies between them, and then
$q + \sqrt 2 \in \intoo{x}{y}$ is irrational (if $q + \sqrt 2$ were
rational, so would be $\sqrt 2$).
\end{proof}

\begin{method}[Proving equalities with sup and inf]\label{met:b1:reals:supproofs}
To prove $\sup A = s$: check that $s$ bounds $A$ above, then produce,
for each $\varepsilon > 0$ (or for a sequence $\varepsilon = \frac
1n$), an element of $A$ above $s - \varepsilon$. To compare suprema,
use: $A \subseteq B \implies \sup A \leq \sup B$; and for all $a, b$:
$\sup(A + B) = \sup A + \sup B$, where $A + B = \{a + b\}$
(\cref{exo:b1:reals:5}). Never write $\sup A$ before knowing $A$ is
nonempty and bounded above.
\end{method}

\section{Intervals}

\begin{proposition}[Characterization of intervals]\label{prop:b1:reals:intervals}
A subset $I \subseteq \R$ is an \emph{interval}\index{interval} (one
of the familiar types $\intoo{a}{b}$, $\intcc{a}{b}$, $\intco{a}{b}$,
$\intoc{a}{b}$, half-lines, $\R$, $\emptyset$, singletons) if and
only if it is \emph{convex}:
\[
\forall x, y \in I,\ \forall z \in \R, \quad
x \leq z \leq y \implies z \in I .
\]
\end{proposition}

\begin{proof}
Every listed type is clearly convex. Conversely, let $I$ be convex and
nonempty. Set $a = \inf I$ if $I$ is bounded below, else $a =
-\infty$; likewise $b = \sup I$ or $+\infty$. We claim $\intoo{a}{b}
\subseteq I \subseteq \intcc{a}{b}$ (with obvious conventions at
$\pm\infty$). The second inclusion is the definition of bounds. For
the first, let $z \in \intoo{a}{b}$: since $z > a$, $z$ is not a
lower bound (or $a = -\infty$), so some $x \in I$ has $x < z$;
similarly some $y \in I$ has $y > z$; convexity puts $z \in I$. Thus
$I$ is $\intoo{a}{b}$ plus possibly its finite endpoints: one of the
listed types.
\end{proof}

\begin{remark}[Extended real line]
It is convenient to adjoin two symbols and work in $\overline\R = \R
\cup \{-\infty, +\infty\}$, with the conventions $\sup A = +\infty$
when $A$ is not bounded above and $\sup \emptyset = -\infty$. Then
\emph{every} subset of $\R$ has a supremum in $\overline\R$ --- a
notational comfort used freely for limits in \cref{ch:b1:seq}.
\end{remark}

\section{Exercises}

\begin{exercise}[$\star$]\label{exo:b1:reals:1}
Determine (with proofs) sup, inf, max, min --- when they exist --- of:
\[
A = \Bigl\{\frac{1}{n} : n \in \N^*\Bigr\},
\qquad
B = \Bigl\{\frac{(-1)^n n}{n+1} : n \in \N\Bigr\},
\qquad
C = \{x \in \R : x^2 < 3\}.
\]
\end{exercise}

\begin{exercise}[$\star$]\label{exo:b1:reals:2}
Prove that for all $x, y \in \R$: $\lfloor x \rfloor + \lfloor y
\rfloor \leq \lfloor x + y \rfloor \leq \lfloor x \rfloor + \lfloor y
\rfloor + 1$, and that both bounds are attained.
\end{exercise}

\begin{exercise}[$\star$]\label{exo:b1:reals:3}
Prove that for every $x \in \R$ and $n \in \N^*$:
$\Bigl\lfloor \frac{\lfloor nx \rfloor}{n} \Bigr\rfloor = \lfloor x
\rfloor$.
\end{exercise}

\begin{exercise}[$\star$]\label{exo:b1:reals:4}
Let $A \subseteq B$ be nonempty subsets of $\R$, $B$ bounded. Prove
$\inf B \leq \inf A \leq \sup A \leq \sup B$.
\end{exercise}

\begin{exercise}[$\star\star$]\label{exo:b1:reals:5}
For nonempty bounded $A, B \subseteq \R$, define $A + B = \{a + b : a
\in A,\ b \in B\}$ and $-A = \{-a : a \in A\}$. Prove:
\[
\sup(A + B) = \sup A + \sup B,
\qquad
\sup(-A) = -\inf A .
\]
\end{exercise}

\begin{exercise}[$\star\star$]\label{exo:b1:reals:6}
Let $f, g \colon E \to \R$ be bounded functions. Prove
\[
\sup_{x \in E}\, \bigl(f(x) + g(x)\bigr) \leq \sup_{x \in E} f(x) +
\sup_{x \in E} g(x),
\]
and give an example where the inequality is strict. Why does this not
contradict \cref{exo:b1:reals:5}?
\end{exercise}

\begin{exercise}[$\star\star$]\label{exo:b1:reals:7}
Prove that $\sqrt 2 + \sqrt 3$ is irrational. \emph{(Square it and
use the irrationality of $\sqrt 6$, to be proved via
\cref{exo:b1:arith:7}.)}
\end{exercise}

\begin{exercise}[$\star\star$]\label{exo:b1:reals:8}
Prove that the set $D = \bigl\{\frac{m}{2^n} : m \in \Z,\ n \in
\N\bigr\}$ of dyadic rationals is dense in $\R$: between any two
reals lies a dyadic rational.
\end{exercise}

\begin{exercise}[$\star\star\star$]\label{exo:b1:reals:9}
Let $G$ be a subgroup of $(\R, +)$ with $G \neq \{0\}$. Set $\alpha =
\inf\,(G \cap \intoo{0}{+\infty})$. Prove:
\begin{enumerate}
  \item if $\alpha > 0$, then $G = \alpha\Z$;
  \item if $\alpha = 0$, then $G$ is dense in $\R$.
\end{enumerate}
Deduce that $\Z + \sqrt 2\,\Z$ is dense in $\R$.
\end{exercise}

\begin{exercise}[$\star\star\star$]\label{exo:b1:reals:10}
For $A, B$ nonempty sets of positive reals, let $AB = \{ab : a \in
A, b \in B\}$. Prove $\sup(AB) = \sup A \cdot \sup B$ (bounded case),
and show by example that positivity is essential.
\end{exercise}
